% =====================================================================
%  Thème 2 — Chasles et combinaisons — Niveau MOYEN — Corrigé
% =====================================================================
\documentclass[11pt,a4paper]{article}
\usepackage[utf8]{inputenc}
\usepackage[T1]{fontenc}
\usepackage{lmodern}
\usepackage[french]{babel}
\usepackage{amsmath,amssymb,mathtools}
\usepackage{xcolor}
\usepackage{enumitem}
\usepackage{fancyhdr}
\usepackage[a4paper,margin=2.1cm,headheight=26pt,headsep=10pt]{geometry}

\definecolor{primary}{RGB}{222,70,80}
\definecolor{primarydark}{RGB}{150,30,42}
\definecolor{excol}{RGB}{0,135,90}

\pagestyle{fancy}\fancyhf{}
\fancyhead[L]{\small\textcolor{primarydark}{\textbf{Fiche 2} — Calcul vectoriel}}
\fancyhead[R]{\small\textcolor{primarydark}{\textsc{Thème 2 — Moyen — Corrigé}}}
\fancyfoot[C]{\small\thepage}
\renewcommand{\headrulewidth}{0.8pt}

\newcommand{\vc}[1]{\overrightarrow{#1}}
\providecommand{\norm}[1]{\left\lVert #1\right\rVert}
\newcommand{\cor}[1]{\medskip\noindent{\color{primarydark}\bfseries\large Corrigé #1.}\par\smallskip}
\newcommand{\rep}{\textcolor{excol}{$\blacktriangleright$}\ }

\begin{document}
\begin{center}
{\LARGE\bfseries\color{primarydark}Chasles et combinaisons — Corrigé Moyen}\\[2pt]
{\color{primary}\rule{0.5\linewidth}{1.1pt}}
\end{center}
\vspace{2mm}

\cor{1}
\begin{enumerate}[label=\textbf{\alph*)},leftmargin=2em,itemsep=2pt]
  \item On réordonne : $\vc{AB}+\vc{BC}+\vc{CD}=\vc{AD}$.
  \item $\vc{AB}-\vc{AC}=\vc{AB}+\vc{CA}=\vc{CA}+\vc{AB}=\vc{CB}$.
  \item $-\vc{DC}=\vc{CD}$, donc $\vc{AB}+\vc{BC}+\vc{CD}=\vc{AD}$.
  \item $\vc{OA}-\vc{OB}=\vc{OA}+\vc{BO}=\vc{BO}+\vc{OA}=\vc{BA}$.
\end{enumerate}
\rep Règle express : $\vc{XA}-\vc{XB}=\vc{BA}$ (origine commune $X$ éliminée). Ex.\ b) et d).

\cor{2}
\begin{enumerate}[label=\textbf{\alph*)},leftmargin=2em,itemsep=2pt]
  \item $\vc{AB}=\vc{AO}+\vc{OB}=\vc{OB}-\vc{OA}$.
  \item $\vc{CA}=\vc{OA}-\vc{OC}$.
  \item $\vc{AB}+\vc{BC}=\vc{AC}=\vc{OC}-\vc{OA}$.
\end{enumerate}
\rep Tout vecteur $\vc{XY}$ s'écrit $\vc{OY}-\vc{OX}$ : c'est le passage aux \emph{vecteurs position}
depuis l'origine $O$.

\cor{3}
\begin{enumerate}[label=\textbf{\alph*)},leftmargin=2em,itemsep=2pt]
  \item $I$ milieu de $[AB]$ : $\vc{AI}=\dfrac12\,\vc{AB}$.
  \item $\vc{IA}+\vc{IB}=\vec0$ : comme $I$ est le milieu, $\vc{IA}$ et $\vc{IB}$ sont opposés
        (même longueur, sens contraires).
  \item $\vc{OI}=\vc{OA}+\vc{AI}=\vc{OA}+\dfrac12\vc{AB}=\vc{OA}+\dfrac12(\vc{OB}-\vc{OA})
        =\dfrac12\vc{OA}+\dfrac12\vc{OB}=\dfrac12(\vc{OA}+\vc{OB})$.
\end{enumerate}
\rep La formule $\vc{OI}=\frac12(\vc{OA}+\vc{OB})$ traduit « le milieu est la moyenne des deux points ».

\cor{4}
\begin{enumerate}[label=\textbf{\alph*)},leftmargin=2em,itemsep=2pt]
  \item $\vc{AC}=\vc{AB}+\vc{AD}=\vec u+\vec v$ (règle du parallélogramme).
  \item $\vc{BD}=\vc{BA}+\vc{AD}=-\vc{AB}+\vc{AD}=\vec v-\vec u$.
\end{enumerate}
\rep Une diagonale est la \emph{somme} des côtés issus du même sommet, l'autre en est la
\emph{différence}.

\cor{5}
\begin{enumerate}[label=\textbf{\alph*)},leftmargin=2em,itemsep=2pt]
  \item $I$ milieu de $[BC]$ : $\vc{BI}=\dfrac12\vc{BC}$. Or $\vc{BC}=\vc{AD}=\vec v$ (côtés opposés),
        donc $\vc{BI}=\dfrac12\vec v$.
  \item $\vc{AI}=\vc{AB}+\vc{BI}=\vec u+\dfrac12\vec v$.
\end{enumerate}
\rep On décompose le trajet $A\to I$ en $A\to B$ puis $B\to I$, puis on remplace chaque morceau par
$\vec u$ ou $\vec v$.

\end{document}
